%%--------------------------------------------------------------------------
%%  SECTION 6: FUTURE RESEARCH DIRECTIONS
%%--------------------------------------------------------------------------
\section{Future Research Directions}

My ongoing and planned research evolves along three interconnected frontiers,
each combining fundamental mathematical questions with high-stakes engineering
applications.

\medskip\noindent\textbf{Foundation models for plasma and fusion confinement.}
Building on the DeepRTE paradigm, I am developing operator-learning
foundation models for the Vlasov--Maxwell system governing fusion plasma
dynamics.  The scientific target is a surrogate model that replaces PIC
simulations for tokamak transport calculations, accurate across the full
kinetic-to-fluid spectrum and evaluable in milliseconds for real-time feedback
control of plasma shape and stability---a prerequisite for commercially viable
fusion energy.  This project is being developed in collaboration with plasma
physicists at SJTU's Institute of Natural Sciences.  A key open challenge is
the preservation of physical invariants (total energy, $L^2$ norm of the
distribution function, entropy inequality) within the operator-learning
architecture, connecting directly to the structure-preserving research thread
below.

\medskip\noindent\textbf{Convergence theory for generative models in inverse problems.}\enspace
Current diffusion-model-based inversion methods, including ODE-DPS, rely on
heuristic justifications for the posterior score approximation step.  I plan
to develop a rigorous error analysis that bounds the Wasserstein distance
between the learned approximate posterior and the true Bayesian posterior as
a function of (i)~the score approximation error, (ii)~the ODE discretization
step size, and (iii)~the forward model accuracy.  This connects F-principle
theory---which governs the spectral convergence of the score network---with
classical Bayesian regularization theory, providing what would be, to our
knowledge, the first end-to-end convergence guarantee for generative inversion
methods.

\medskip\noindent\textbf{Structure-preserving operator learning.}\enspace
Inspired by symplectic integrators and AP schemes in classical numerical
analysis, I aim to build operator-learning architectures that provably
preserve nonlinear conservation laws ($L^p$ mass, energy, entropy inequality)
for Hamiltonian and dissipative PDE systems.  The central technical challenge
is to enforce these constraints as hard architectural invariants rather than
loss-function penalties, since penalty-based approaches are known to degrade
in the stiff multiscale limit---precisely the regime of greatest scientific
interest.  Success here would generalize the APNN design principle to a broad
class of physical constraints applicable to fluid mechanics, plasma physics,
and molecular dynamics simulation.
